\chapter{Powierzchnie algebraiczne: przecięcia, rozdmuchania i szczególne klasy}
\label{ch:ga6-powierzchnie}

\section{Dlaczego geometria powierzchni jest jakościowo nowa?}
\label{sec:ga6-motywacja}

W przypadku gładkiej krzywej algebraicznej punkty są dywizorami,
a odwzorowanie biwymierne między gładkimi krzywymi właściwymi jest
izomorfizmem. W wymiarze dwa sytuacja zmienia się zasadniczo:
można rozdmuchać punkt gładkiej powierzchni i otrzymać inną gładką
powierzchnię. Nowa powierzchnia ma nawet nową krzywą $E\cong\PP^1$,
ale jej ciało funkcji jest takie samo. Żeby odróżnić powierzchnie,
nie wystarcza zatem samo ciało funkcji ani zbiór punktów.

W tym rozdziale głównym obiektem będzie \emph{gładka, spójna,
rzutowa powierzchnia} $S$ nad ciałem algebraicznie domkniętym
$k$ charakterystyki $0$. Oznacza to całkowity schemat gładki i
rzutowy nad $k$, o wymiarze dwa. Gdy nie korzystamy z rzutowości,
powiemy to jawnie. Przypomnijmy, że \emph{całkowity} znaczy
nierozkładalny i zredukowany. W celu porównania z topologią
przyjmujemy $k=\C$.

\begin{przyklad}[Pierwsze cztery powierzchnie]
\label{prz:ga6-examples}
\begin{enumerate}[label=(\roman*)]
\item $\PP^2_k$ jest gładką powierzchnią rzutową; proste w niej
      mają wspólny stopień jeden.
\item $\PP^1_k\times\PP^1_k$ jest gładką powierzchnią z dwiema
      różnymi rodzinami krzywych będących włóknami rzutowań.
\item Gładka powierzchnia kubiczna $S=\V_+(F)\subseteq\PP^3_k$,
      gdzie $F$ jest jednorodny stopnia trzy, zawiera dokładnie
      $27$ prostych (udowodnimy część rachunkową tego wyniku).
\item Gładka powierzchnia stopnia czwartego w $\PP^3_\C$
      jest powierzchnią K3; zobaczymy, dlaczego jej wiązka
      kanoniczna jest trywialna i dlaczego nie jest wymierna.
\end{enumerate}
\end{przyklad}

\begin{definicja}[Wymierność i powierzchnia wymierna]
Niech $X,Y$ będą całkowitymi rozmaitościami algebraicznymi nad $k$.
Są \emph{biwymiernie równoważne}, gdy zawierają izomorficzne
niepuste otwarte podzbiory; równoważnie $k(X)\cong k(Y)$ jako
rozszerzenia ciała $k$. Powierzchnia $S$ jest \emph{wymierna},
jeżeli jest biwymiernie równoważna z $\PP^2_k$.
\end{definicja}

\begin{uwaga}[Dwa rodzaje klasyfikacji]
Klasyfikacja \emph{z dokładnością do izomorfizmu} odróżnia
$\PP^2$ od jego rozdmuchania w punkcie. Klasyfikacja
\emph{biwymierna} je utożsamia. Jeżeli chcemy zachować
liczby przecięć i układ krzywych, musimy kontrolować
konkretne modele, nie tylko ich ciało funkcji.
\end{uwaga}

\section{Przecięcia krzywych: od prostych rachunków do formy przecięć}
\label{sec:ga6-intersections}

\subsection{Krotność lokalna: dlaczego samo liczenie punktów nie działa?}

Niech $S$ będzie gładką powierzchnią nad $k$ oraz $C,D\subseteq S$
dwiema \emph{różnymi} całkowitymi krzywymi. W otoczeniu danego
punktu $p\in C\cap D$ krzywe są zadane przez pojedyncze równania
lokalne $f=0$ i $g=0$. Istnienie takich równań wynika z faktu,
że regularny pierścień lokalny powierzchni jest pierścieniem
z jednoznacznym rozkładem na czynniki. To ostatnie twierdzenie
o regularnych pierścieniach lokalnych przyjmujemy na tym etapie
jako wynik algebry przemiennej.

\begin{definicja}[Lokalna krotność przecięcia]
\label{def:ga6-intersection-local}
Jeżeli $p\in C\cap D$ jest izolowanym punktem ich przecięcia,
a $f,g\in\mathcal O_{S,p}$ są lokalnymi równaniami $C,D$,
to definiujemy
\[
 i_p(C,D;S):=\operatorname{length}_{\mathcal O_{S,p}}
   \bigl(\mathcal O_{S,p}/(f,g)\bigr).
\]
Długość modułu oznacza liczbę czynników w jego szeregu
kompozycyjnym; dla $k$ algebraicznie domkniętego, przy
skończonej długości, jest równa wymiarowi tej przestrzeni
jako przestrzeni wektorowej nad $k$. Krzywe przecinają się
\emph{transwersalnie} w $p$, gdy ich przestrzenie styczne w $p$
są różnymi prostymi w $T_pS$.
\end{definicja}

\begin{stwierdzenie}[Niezależność od wyboru równań]
\label{st:ga6-independent-local}
Niech $S$ będzie gładką powierzchnią nad ciałem $k$,
a $C,D\subseteq S$ efektywnymi dzielnikami Cartiera
bez wspólnej składowej w otoczeniu punktu zamkniętego
$p\in C\cap D$. Długość
$\operatorname{length}_{\mathcal O_{S,p}}
\bigl(\mathcal O_{S,p}/(f,g)\bigr)$, gdzie $f,g$
są lokalnymi równaniami $C,D$, nie zależy od ich wyboru.
\end{stwierdzenie}
\begin{proof}
Dwa równania definiujące ten sam efektywny dzielnik
Cartiera w otoczeniu $p$ różnią się o jednostkę:
$f'=uf$ i $g'=vg$, gdzie $u,v\in\mathcal O_{S,p}^{\times}$.
Zatem $(f',g')=(f,g)$ i pierścienie ilorazowe są takie same.
Wobec tego długość nie zależy od lokalnych generatorów.
\end{proof}

\begin{stwierdzenie}[Przecięcie transwersalne ma krotność jeden]
\label{st:ga6-transversal-one}
Niech $p$ będzie punktem gładkim obu krzywych $C,D$ na
gładkiej powierzchni $S$. Wówczas $i_p(C,D;S)=1$ wtedy
i tylko wtedy, gdy krzywe przecinają się transwersalnie w $p$.
\end{stwierdzenie}
\begin{proof}
Oznaczmy $R=\mathcal O_{S,p}$ i $\mathfrak m=\mathfrak m_p$.
Ponieważ $R$ jest regularnym pierścieniem lokalnym wymiaru dwa,
przestrzeń $\mathfrak m/\mathfrak m^2$ ma wymiar dwa nad $k$.
Różniczki $df(p),dg(p)$ są niezależne wtedy i tylko wtedy,
gdy klasy $f,g$ tworzą bazę $\mathfrak m/\mathfrak m^2$.
Z lematu Nakayamy otrzymujemy wówczas $(f,g)=\mathfrak m$,
więc $R/(f,g)=R/\mathfrak m\cong k$ ma długość jeden.
Odwrotnie, jeśli długość $R/(f,g)$ jest równa jeden,
to iloraz jest prostym $R$-modułem; ponieważ $(f,g)$ jest
$\mathfrak m$-pierwotny, musi zachodzić $(f,g)=\mathfrak m$.
Klasy $f,g$ generują wtedy dwuwymiarową przestrzeń
$\mathfrak m/\mathfrak m^2$ i są liniowo niezależne.
Jest to właśnie warunek transwersalności.
\end{proof}

\begin{przyklad}[Dwie krzywe styczne]
\label{prz:ga6-tangent}
Na $\A^2_{x,y}$ rozważmy $C=\V(y)$ i $D=\V(y-x^m)$,
gdzie $m\ge1$. Obie krzywe są gładkie w początku.
Ponieważ w pierścieniu lokalnym można najpierw podstawić
$y=0$, mamy
$\mathcal O_{\A^2,0}/(y,y-x^m)\cong k[x]_{(x)}/(x^m)$.
Jego bazę nad $k$ stanowią klasy $1,x,\ldots,x^{m-1}$;
więc $i_0(C,D;\A^2)=m$. Dla $m=1$ przecięcie jest
transwersalne, dla $m=2$ jest styczne pierwszego rzędu,
a dowolnie wysokie $m$ daje dowolnie dużą krotność mimo
obecności tylko jednego punktu przecięcia.
\end{przyklad}

\begin{center}
\begin{minipage}{\linewidth}
\centering
\begin{tikzpicture}[scale=0.88,>=Stealth,every node/.style={font=\small}]
  \begin{scope}[shift={(-3.2,0)}]
   \draw[gray!65,->] (-1.45,0)--(1.45,0) node[below right] {$x$};
   \draw[gray!65,->] (0,-0.9)--(0,1.45) node[left] {$y$};
   \draw[blue!70!black,thick] (-1.05,-0.7)--(1.05,0.7);
   \draw[red!70!black,thick] (-1.05,0.7)--(1.05,-0.7);
   \node[align=center,text width=4.3cm] at (0,-1.65)
     {Przecięcie transwersalne\\$i_0=1$};
  \end{scope}
  \begin{scope}[shift={(3.2,0)}]
   \draw[gray!65,->] (-1.45,0)--(1.45,0) node[below right] {$x$};
   \draw[gray!65,->] (0,-0.9)--(0,1.45) node[left] {$y$};
   \draw[red!70!black,thick] (-1.3,0)--(1.3,0);
   \draw[blue!70!black,thick,domain=-1.1:1.1,variable=\t]
    plot ({\t},{0.8*\t*\t});
   \node[align=center,text width=4.3cm] at (0,-1.65)
      {Przecięcie styczne\\$i_0=2$};
  \end{scope}
\end{tikzpicture}

{\small Rysunek 6.1. Jeden punkt przecięcia może mieć różną krotność algebraiczną. W lewym panelu różniczki równań są niezależne; w prawym krzywe mają wspólną styczną.}\par
\end{minipage}
\end{center}

\subsection{Globalne liczby przecięć i samoprzecięcie}

\begin{definicja}[Dywizory i ich równoważność liniowa]
\label{def:ga6-divisor-group}
Niech $S$ będzie gładką całkowitą powierzchnią nad ciałem $k$.
Grupa $\operatorname{Div}(S)$ jest wolną grupą abelową na
całkowitych domkniętych krzywych $C\subseteq S$. Jej elementy
$D=\sum_C n_C C$ są \emph{dywizorami Weila}; suma ma skończony
nośnik. Dla niezerowej funkcji wymiernej $f\in k(S)^\times$
definiujemy $\operatorname{div}(f)=\sum_C\operatorname{ord}_C(f)C$.
Rząd $\operatorname{ord}_C$ jest wartościowaniem w pierścieniu lokalnym
$\mathcal O_{S,\eta_C}$ punktu generycznego krzywej $C$;
jest to dyskretny pierścień wartościowania, ponieważ $S$ jest
regularna. Dwa dywizory $D,D'$ są \emph{liniowo równoważne},
co oznaczamy $D\sim D'$, gdy $D-D'=\operatorname{div}(f)$
dla pewnego $f\in k(S)^\times$.
\end{definicja}

\begin{twierdzenie}[Forma przecięć na powierzchni rzutowej]
\label{tw:ga6-intersection-global}
Niech $S$ będzie gładką rzutową powierzchnią nad $k$.
Istnieje jednoznaczna symetryczna forma dwuliniowa
$\operatorname{Div}(S)\times\operatorname{Div}(S)\to\Z$,
$(D_1,D_2)\mapsto D_1\cdot D_2$, zależna tylko od klas
równoważności liniowej, która spełnia:
\begin{enumerate}[label=(\roman*)]
\item dla różnych krzywych całkowitych $C,D$ zachodzi
$C\cdot D=\sum_{p\in C\cap D}i_p(C,D;S)$;
\item dla gładkiej krzywej całkowitej $C\subseteq S$
zachodzi $C^2:=C\cdot C=\deg\bigl(\mathcal O_S(C)|_C\bigr)$;
\item dla dowolnej krzywej całkowitej $C$ i dywizora $D$
zachodzi $D\cdot C=\deg(\mathcal O_S(D)|_{\widetilde C})$,
gdzie $\widetilde C\to C$ jest normalizacją.
\end{enumerate}
W (iii) stopień wiązki na normalizacji liczymy z jej
pullbacku. Na gładkiej krzywej $C$ wzór upraszcza się
do stopnia na samej $C$.
\end{twierdzenie}

\begin{uwaga}[Granica pełnego dowodu]
Nie będziemy udawać, że globalna postać formy przecięć
wynika wyłącznie z dwóch przykładów lokalnych. Jej konstrukcja
wymaga m.in. zgodności lokalnej długości z teorią klas dywizorów
i stopni wiązek oraz lematu o przesuwaniu cykli.
\emph{Twierdzenie~\ref{tw:ga6-intersection-global} przyjmujemy tutaj
jako podstawowy wynik teorii przecięć}. Jego lokalne części
są wyliczone bezpośrednio, natomiast kluczowe zastosowania
do rozdmuchań poniżej zostaną uzasadnione rachunkiem wiązek.
W rozdziale 7 wyprowadzimy formę przecięć przez
pierwszą klasę Cherna i stopnie na krzywych;
sam ogólny lemat o przesuwaniu cykli zostanie tam jawnie
odnotowany jako osobne twierdzenie.
\end{uwaga}

\begin{przyklad}[Przecięcia w $\PP^2$]
\label{prz:ga6-p2}
Niech $H\subseteq\PP^2$ będzie prostą. Dwie różne proste
przecinają się w jednym punkcie transwersalnie,
więc $H^2=1$ (zależność od klasy liniowej pozwala zastąpić
prostą inną prostą). Krzywa płaska zadana równaniem jednorodnym
stopnia $d$ ma klasę liniową $dH$: jest to dywizor zer
sekcji $\mathcal O_{\PP^2}(d)$. Z dwuliniowości wynika
$(dH)\cdot(eH)=de$, czyli wersja twierdzenia Bézouta
z krotnościami, jeśli dwie krzywe nie mają wspólnej składowej.
Dla krzywej stopnia $d$ jej \emph{samoprzecięcie} wynosi $d^2$,
choć geometrycznie nie przecinamy krzywej z samą sobą
w sensie zbiorowym.
\end{przyklad}

\begin{definicja}[Wiązka normalna krzywej]
Jeżeli $C\subseteq S$ jest gładką całkowitą krzywą
na gładkiej powierzchni, to jej \emph{wiązka normalna}
jest wiązką liniową
$N_{C/S}:=(\mathcal I_C/\mathcal I_C^2)^\vee$ na $C$.
Z identyfikacji snopa ideałów $\mathcal I_C\cong\mathcal O_S(-C)$
otrzymujemy $N_{C/S}\cong\mathcal O_S(C)|_C$, a zatem
$C^2=\deg N_{C/S}$. To nadaje samoprzecięciu sens
geometryczny: mierzy stopień normalnego kierunku do krzywej.
\end{definicja}

\section{Rozdmuchanie punktu: postać przecięć i klasa kanoniczna}
\label{sec:ga6-blowup}

Niech teraz $\pi:\widetilde S\to S$ będzie rozdmuchaniem
gładkiej rzutowej powierzchni w punkcie $p\in S$,
a $E=\pi^{-1}(p)\cong\PP^1$ jego dzielnikiem wyjątkowym.

% Wspólna podstawa geometryczna dla rozdmuchania punktu.
% Wariant opisu przez normalną i przez konormalną: obydwa zachowane.
% Historyczne etykiety występują tylko przy pierwszym włączeniu.
\ifcsname biblioteka@rozdmuchanie@normalna@loaded\endcsname
Wiązkę normalną dzielnika wyjątkowego obliczono wcześniej
w lemacie~\ref{lem:xvii-exceptional}.
\else
\expandafter\gdef\csname biblioteka@rozdmuchanie@normalna@loaded\endcsname{}
\begin{lemat}[Wiązka normalna dzielnika wyjątkowego]
\label{st:ga6-exceptional-normal}\label{lem:xvii-exceptional}
Niech $S$ będzie gładką zespoloną powierzchnią, $p\in S$, a
$\pi:\widetilde S=\operatorname{Bl}_pS\to S$ jej rozdmuchaniem
w punkcie $p$. Wtedy $E=\pi^{-1}(p)\simeq\mathbb P^1_{\mathbb C}$,
$N_{E/\widetilde S}\simeq\mathcal O_{\mathbb P^1}(-1)$ i
$E\cdot E=\deg N_{E/\widetilde S}=-1$.
Ostatnią liczbę rozumiemy jako lokalne samoprzecięcie krzywej,
nawet jeśli $S$ nie jest zwarta.
\end{lemat}
\begin{proof}
Wybierzmy współrzędne $(x,y)$ powierzchni $S$ z początkiem w $p$.
Rozdmuchanie jest lokalnie pododmianą
$\{((x,y),[s:t])\colon xt=ys\}\subset\mathbb C^2\times\mathbb P^1$.
W karcie $s\ne0$ piszemy $v=t/s$, $y=xv$, a w karcie $t\ne0$
piszemy $u=s/t$, $x=uy$. Na części wspólnej zachodzi
$u=v^{-1}$ i $y=vx$. W pierwszej karcie ideał $E$ jest
generowany przez $x$, w drugiej przez $y$.

\emph{Opis przez wiązkę normalną.}
W punkcie $[1:v]\in E$ kierunek normalny do $E$ jest
reprezentowany przez wektor $(x,xv)=x(1,v)$.
Po zmianie karty wektory te wyznaczają dokładnie
podwiązki prostych rozpiętych przez $(s,t)$ w
$\mathbb P^1\times\mathbb C^2$; jest to tautologiczna
wiązka $\mathcal O_{\mathbb P^1}(-1)$.

\emph{Równoważny opis przez wiązkę konormalną.}
Lokalne generatory $x$ i $y$ wiązki
$\mathcal I_E/\mathcal I_E^2$ spełniają $y=vx$;
odpowiada to wiązce $\mathcal O_{\mathbb P^1}(1)$.
Jej dualna jest więc $N_{E/\widetilde S}\cong\mathcal O(-1)$.
Ponieważ stopień $\mathcal O_{\mathbb P^1}(-1)$ wynosi $-1$,
formuła samoprzecięcia
$E^2=\langle c_1(N_{E/\widetilde S}),[E]\rangle$
daje $E^2=-1$.
\end{proof}
\fi


\begin{stwierdzenie}[Ortogonalność i formuła przecięć po rozdmuchaniu]
\label{st:ga6-blowup-form}
Niech $D,D'$ będą dywizorami na $S$. Wtedy
\[
 (\pi^*D)\cdot E=0,\qquad
 (\pi^*D)\cdot(\pi^*D')=D\cdot D',\qquad E^2=-1.
\]
Jeżeli całkowita krzywa $C\subseteq S$ przechodzi
przez $p$ z krotnością $m$, to jej transformata ścisła
$\widetilde C$ spełnia
$\pi^*C=\widetilde C+mE$ jako dywizory, a zatem
\[
 \widetilde C^2=C^2-m^2,\qquad
 \widetilde C\cdot E=m.
\]
\end{stwierdzenie}
\begin{proof}
Ponieważ $\pi|_E:E\to\{p\}$ jest odwzorowaniem stałym,
$\mathcal O_{\widetilde S}(\pi^*D)|_E$
jest trywialną wiązką $\mathcal O_E$;
jej stopień wynosi zero. Z definicji formy przecięć
otrzymujemy $(\pi^*D)\cdot E=0$.
Druga równość jest szczególną postacią \emph{wzoru projekcji}
dla przecięć: $\pi_*(\pi^*D\cdot\pi^*D')=D\cdot D'$.
W przypadku $\PP^2$ można ją też zobaczyć bez ogólnego
wzoru projekcji: przesuwamy przedstawicieli klas
prostych poza $p$, gdzie $\pi$ jest izomorfizmem.
Ogólną wersję wzoru projekcji traktujemy jako część
przyjętej teorii przecięć.

Dla ostatnich równości weźmy w $p$ lokalne równanie
$f$ krzywej $C$. Jeśli jego najmniejsza część
jednorodna ma stopień $m$, to na charcie $y=xv$
znany z rozdziału 5 wzór daje
$f(x,xv)=x^mF(x,v)$. Wyrażenie $x^m$ oznacza $mE$,
a $F=0$ transformatę ścisłą. Analogiczny wzór na drugim
charcie skleja się do równości dywizorów
$\pi^*C=\widetilde C+mE$. Stąd, korzystając z
$\pi^*C\cdot E=0$ oraz $E^2=-1$, mamy
$\widetilde C\cdot E=m$. Wreszcie
$\widetilde C^2=(\pi^*C-mE)^2=C^2-m^2$.
\end{proof}

\begin{przyklad}[Zerowanie i zmiana samoprzecięcia]
\label{prz:ga6-blowup-curves}
W $S=\PP^2$ prosta ma klasę $H$ i $H^2=1$.
Po rozdmuchaniu punktu $p$ należącego do prostej
jej transformata ścisła ma klasę $H-E$ i
$(H-E)^2=1-1=0$. Po rozdmuchaniu \emph{dwóch różnych}
punktów tej samej prostej otrzymujemy klasę
$H-E_1-E_2$ o samoprzecięciu $-1$.
Jeżeli rozdmuchiwany punkt nie leży na prostej,
transformata ma klasę $H$ i samoprzecięcie nadal $1$.
Ta sama krzywa jako abstrakcyjna $\PP^1$ może więc
mieć w różnych powierzchniach wiązki normalne
różnych stopni. Samoprzecięcie nie jest niezmiennikiem
samej krzywej, lecz jej zanurzenia.
\end{przyklad}

\begin{stwierdzenie}[Klasa kanoniczna po rozdmuchaniu]
\label{st:ga6-Kblowup}
Niech $K_S$ oznacza dowolny dywizor kanoniczny
gładkiej powierzchni $S$. Wtedy
\[
 K_{\widetilde S}\sim\pi^*K_S+E,
 \qquad K_{\widetilde S}^2=K_S^2-1.
\]
Pierwsza równość jest równością klas liniowej równoważności.
\end{stwierdzenie}
\begin{proof}
W lokalnych współrzędnych $x,y$ na $S$ generator
wiązki kanonicznej jest reprezentowany przez
$dx\wedge dy$. Na charcie rozdmuchania $y=xv$
mamy $dy=v\,dx+x\,dv$, zatem
$\pi^*(dx\wedge dy)=x\,dx\wedge dv$.
Przy przeciąganiu lokalnej różniczki dwuwymiarowej
pojawia się dokładnie jeden czynnik $x$ — równanie
$E$ na tym charcie. Na drugim charcie analogicznie
$dx\wedge dy=-(y)\,du\wedge dy$ po cofnięciu,
więc również współczynnik zanika na $E$ rzędu jeden.
W ten sposób dywizor kanoniczny na $\widetilde S$
równa się klasie $\pi^*K_S+E$.
Równość kwadratów wynika z poprzedniego stwierdzenia:
$(\pi^*K_S+E)^2=K_S^2+E^2=K_S^2-1$.
\end{proof}

\begin{stwierdzenie}[Grupa Picarda po rozdmuchaniu]
\label{st:ga6-pic-blowup}
Dla $\pi:\widetilde S=\operatorname{Bl}_pS\to S$,
gdzie $S$ jest gładką rzutową powierzchnią, zachodzi
\[
 \operatorname{Pic}(\widetilde S)\cong
 \pi^*\operatorname{Pic}(S)\oplus\Z[E].
\]
W szczególności dla rozdmuchania $\PP^2$ w $r$
różnych punktach otrzymujemy wolną grupę abelową
z bazą $H,E_1,\ldots,E_r$; w tej bazie formę
przecięć opisuje diagonalna macierz
$\operatorname{diag}(1,-1,\ldots,-1)$.
\end{stwierdzenie}
\begin{proof}
Na gładkiej powierzchni każdy dzielnik Weila jest
Cartiera: wykorzystujemy standardowe twierdzenie,
że regularne pierścienie lokalne są UFD. Niech
$\widetilde D$ będzie dzielnikiem na $\widetilde S$.
Każdą jego nierozkładalną składową różną od $E$
można zrzutować na krzywą $D\subseteq S$; jej
transformata ścisła różni się od $\pi^*D$ o
całkowitą wielokrotność $E$. Po zsumowaniu składowych
otrzymujemy $\widetilde D\sim\pi^*D+aE$.
To dowodzi surjektywności odpowiedniego homomorfizmu.

Jeżeli $\pi^*D+aE$ jest liniowo równoważny zeru,
to po ograniczeniu wiązki $\mathcal O_{\widetilde S}(\pi^*D+aE)$
do $E$ otrzymujemy $\mathcal O_{\PP^1}(-a)$.
Musi ona być trywialna, więc $a=0$. Wtedy
$\pi^*D\sim0$. Zmniejszając powierzchnie do
dopełnień punktu $p$ i krzywej $E$, na których
$\pi$ jest izomorfizmem, otrzymujemy $D|_{S\setminus\{p\}}\sim0$.
Na gładkiej powierzchni usunięcie punktu kowymiaru
dwa nie zmienia klas dywizorów: przedstawiciele
Weila mają nośniki kowymiaru jeden, a funkcje
wymierne są te same. Wobec tego $D\sim0$.
Uzyskaliśmy również injektywność i rozkład bezpośredni.
Dla $S=\PP^2$ jej grupa Picarda jest generowana
przez $H=\mathcal O_{\PP^2}(1)$; kolejne rozdmuchania
powtarzają dowód i dają podaną bazę.
\end{proof}

\begin{uwaga}[Kiedy klasa $E_1$ nie oznacza już tej samej krzywej?]
\label{uw:ga6-infinitely-near}
Jeżeli po rozdmuchaniu $p$ rozdmuchujemy następnie punkt
$q$ \emph{na pierwszym dzielniku wyjątkowym}, to
$E_1$ w bazie grupy Picarda oznacza \emph{całkowity
przeciwobraz} poprzedniego dzielnika, nie jego aktualną
transformatę ścisłą. Ta ostatnia ma klasę $E_1-E_2$,
a więc samoprzecięcie $(E_1-E_2)^2=-2$.
Ponadto $(E_1-E_2)\cdot E_2=1$. Otrzymujemy dwie
gładkie krzywe przecinające się transwersalnie,
z samoprzecięciami $-2$ i $-1$. To właśnie typ
konfiguracji, który pojawia się przy rozwiązywaniu
osobliwości.
\end{uwaga}

\begin{center}
\begin{minipage}{\linewidth}
\centering
\begin{tikzpicture}[>=Stealth,every node/.style={font=\small}]
 \node[align=center,draw,rounded corners,inner sep=8pt,text width=4.1cm] (a) at (0,0)
 {Rozdmuchanie pierwszego punktu\\$E_1^2=-1$};
 \node[align=center,draw,rounded corners,inner sep=8pt,text width=4.6cm] (b) at (7.5,0)
 {Drugie rozdmuchanie na $E_1$\\$(E_1-E_2)^2=-2$, $E_2^2=-1$};
 \draw[->,thick] (a.east)--node[above] {punkt na $E_1$}(b.west);
 \draw[red!70!black,very thick] (-0.65,-1.55)--(0.65,-1.55);
 \node[below] at (0,-1.65) {$E_1$};
 \draw[red!70!black,very thick] (6.7,-1.55)--(7.85,-1.55);
 \draw[blue!65!black,very thick] (7.85,-2.05)--(7.85,-1.05);
 \fill (7.85,-1.55) circle (2pt);
 \node[below] at (6.95,-1.65) {$E_1'$};
 \node[right] at (7.95,-1.2) {$E_2$};
\end{tikzpicture}

{\small Rysunek 6.2. Po drugim rozdmuchaniu zmienia się samoprzecięcie transformaty ścisłej starego dzielnika. Nie należy utożsamiać jej klasy z klasą całkowitego przeciwobrazu.}\par
\end{minipage}
\end{center}

\section{Wzór adjunkcji i krzywe wyjątkowe}
\label{sec:ga6-adjunction}

\begin{twierdzenie}[Wzór adjunkcji dla gładkiej krzywej na powierzchni]
\label{tw:ga6-adjunction}
Niech $S$ będzie gładką rzutową powierzchnią, a
$C\subseteq S$ gładką, spójną krzywą. Wówczas
$K_C\sim (K_S+C)|_C$, a zatem
\[
 2g(C)-2=(K_S+C)\cdot C=K_S\cdot C+C^2.
\]
\end{twierdzenie}
\begin{proof}
Ponieważ $C$ jest gładkim efektywnym dzielnikiem
Cartiera, otrzymujemy dokładny ciąg snopów na $C$
\[
0\longrightarrow\mathcal I_C/\mathcal I_C^2
 \longrightarrow\Omega^1_S|_C\longrightarrow\Omega^1_C
 \longrightarrow0.
\]
Pierwszy wyraz jest wiązką liniową
$\mathcal O_C(-C)$; drugi ma rząd dwa, a trzeci
rząd jeden. Biorąc wyznaczniki tego dokładnego
ciągu wiązek lokalnie wolnych, otrzymujemy
$\det(\Omega^1_S|_C)\cong
 \mathcal O_C(-C)\otimes\Omega^1_C$.
Lewa strona jest $\mathcal O_C(K_S)$,
a więc po przeniesieniu czynnika
$\mathcal O_C(-C)$ otrzymujemy
$\Omega_C^1\cong\mathcal O_C(K_S+C)$.
To daje pierwszą równość. Druga wynika przez
wzięcie stopni oraz ze znanego z twierdzenia
Riemanna--Roch faktu $\deg K_C=2g(C)-2$.
\end{proof}

\begin{wniosek}[Samoprzecięcie wymiernej krzywej]
\label{wn:ga6-rational-adjunction}
Jeżeli $C\cong\PP^1$ jest gładką krzywą na gładkiej
rzutowej powierzchni $S$, to
$C^2=-2-K_S\cdot C$.
W szczególności jeżeli $K_S\cdot C=-1$, to $C^2=-1$.
\end{wniosek}
\begin{proof}
Wystarczy wstawić $g(C)=0$ do wzoru adjunkcji.
\end{proof}

\begin{definicja}[Krzywa wyjątkowa pierwszego rodzaju]
Gładką krzywą $E\subseteq S$, izomorficzną z
$\PP^1$ i spełniającą $E^2=-1$, nazywamy
\emph{krzywą $(-1)$}, czyli krzywą wyjątkową
pierwszego rodzaju. Rozdmuchanie punktu na gładkiej
powierzchni wytwarza taką krzywą.
\end{definicja}

\begin{twierdzenie}[Castelnuovo: kryterium zdmuchnięcia]
\label{tw:ga6-castelnuovo}
Niech $S$ będzie gładką rzutową powierzchnią nad
algebraicznie domkniętym ciałem charakterystyki zero.
Jeśli $E\subseteq S$ jest krzywą $(-1)$, to istnieje
gładka powierzchnia $S'$ i właściwy morfizm
$\rho:S\to S'$ kontraktujący dokładnie $E$ do jednego
gładkiego punktu i będący izomorfizmem poza $E$.
Morfizm $\rho$ jest odwrotnością rozdmuchania tego punktu.
\end{twierdzenie}

\begin{uwaga}[Status dowodu Castelnuova]
Nie wyprowadzamy tego twierdzenia z samego warunku
$E^2=-1$: wymaga ono argumentu o deformacjach,
wiązce normalnej i budowie lokalnego otoczenia krzywej.
Widzieliśmy natomiast \emph{model odwrotny}: wyjątkowy
dzielnik rozdmuchania gładkiego punktu ma właśnie
wiązkę normalną $\mathcal O(-1)$. Twierdzenie
Castelnuova mówi, że w kategorii gładkich powierzchni
rzutowych jest to jedyny możliwy lokalny model
gładkiego zdmuchnięcia. Pozostaje jawnie przyjętym
wynikiem zewnętrznym.
\end{uwaga}

\begin{przyklad}[Krzywa $(-2)$ nie może zostać zdmuchnięta do gładkiego punktu]
Rozważmy dwa kolejne rozdmuchania z
uwagi~\ref{uw:ga6-infinitely-near}. Transformata
ścisła pierwszego dzielnika ma klasę $E_1-E_2$
i kwadrat $-2$. Nie spełnia więc warunku
Castelnuova, choć nadal jest krzywą $\PP^1$.
Może być kontraktowana w innym kontekście do
punktu osobliwego, lecz nie jest obrazem wyjątkowej
krzywej jednego rozdmuchania \emph{gładkiego} punktu.
\end{przyklad}

\section{Powierzchnie wymierne i powierzchnie Hirzebrucha}
\label{sec:ga6-rational}

\begin{przyklad}[Dwie rodziny prostych na $\PP^1\times\PP^1$]
Niech $S=\PP^1\times\PP^1$ i niech
$F_1=\{p\}\times\PP^1$, $F_2=\PP^1\times\{q\}$.
Dwa różne włókna tego samego rzutowania są rozłączne,
więc $F_1^2=F_2^2=0$ przez równoważność liniową.
Natomiast $F_1$ i $F_2$ przecinają się dokładnie
w jednym punkcie transwersalnie, zatem
$F_1\cdot F_2=1$. Dla dywizora $aF_1+bF_2$
mamy więc $(aF_1+bF_2)^2=2ab$.
Kanoniczny dywizor jest równy
$K_S\sim-2F_1-2F_2$, ponieważ
$K_{\PP^1\times\PP^1}$ jest tensorowym iloczynem
pullbacków wiązek $K_{\PP^1}=\mathcal O(-2)$.
Stąd $K_S^2=8$.
\end{przyklad}

\begin{definicja}[Powierzchnia Hirzebrucha]
Dla liczby całkowitej $n\ge0$ \emph{powierzchnią
Hirzebrucha} $\mathbb F_n$ nazywamy
$\PP_{\PP^1}\bigl(\mathcal O_{\PP^1}\oplus
\mathcal O_{\PP^1}(n)\bigr)$, z konwencją, w której
$\PP$ parametryzuje jednowymiarowe \emph{ilorazy}
włókien wiązki wektorowej. Ma ona naturalne rzutowanie
$f:\mathbb F_n\to\PP^1$, którego włókna są prostymi
rzutowymi. Oznaczymy przez $F$ klasę włókna,
a przez $B$ sekcję odpowiadającą ilorazowi
$\mathcal O\oplus\mathcal O(n)\twoheadrightarrow\mathcal O$.
\end{definicja}

\begin{stwierdzenie}[Przecięcia na powierzchni Hirzebrucha]
\label{st:ga6-hirzebruch}
W $\mathbb F_n$ zachodzą równości
\[
 B^2=-n,\qquad B\cdot F=1,\qquad F^2=0,
 \qquad K_{\mathbb F_n}\sim-2B-(n+2)F.
\]
W szczególności $K_{\mathbb F_n}^2=8$ dla każdego $n\ge0$.
\end{stwierdzenie}
\begin{proof}
Rzutowanie $f$ ma rozłączne włókna liniowo równoważne,
więc $F^2=0$. Sekcja przecina każde włókno w jednym
punkcie transwersalnie, stąd $B\cdot F=1$.
Wiązka normalna sekcji określonej przez iloraz
$\mathcal O\oplus\mathcal O(n)\twoheadrightarrow\mathcal O$
jest homomorfizmami z jego jądra $\mathcal O(n)$
do ilorazu $\mathcal O$, a więc
$N_{B/\mathbb F_n}\cong\operatorname{Hom}
(\mathcal O(n),\mathcal O)\cong\mathcal O(-n)$.
Zatem $B^2=\deg N_{B/\mathbb F_n}=-n$.

Grupa Picarda $\mathbb F_n$ jest generowana przez
$B,F$ (to standardowa własność projektywizacji
wiązki; można ją zobaczyć z lokalnych trywializacji
i funkcji przejścia). Zapiszmy
$K_{\mathbb F_n}\sim aB+bF$. Adjunkcja na włóknie
$F\cong\PP^1$ daje $K\cdot F=-2-F^2=-2$,
czyli $a=-2$. Adjunkcja na sekcji $B\cong\PP^1$
daje $K\cdot B=-2-B^2=n-2$.
Zatem $-2(-n)+b=n-2$, skąd $b=-n-2$.
Podstawienie do formy przecięć daje
$K^2=(-2B-(n+2)F)^2=4(-n)+4(n+2)=8$.
\end{proof}

\begin{przyklad}[Dlaczego $\mathbb F_1$ jest rozdmuchaniem płaszczyzny?]
\label{prz:ga6-f1}
Ustalmy $p=[0:0:1]\in\PP^2$ i współrzędne $[x:y:z]$.
Rozdmuchanie można zapisać jako powierzchnię incydencji
\[
B=\{([x:y:z],[u:v])\in\PP^2\times\PP^1:\ xv=yu\}.
\]
Projekcja na $\PP^2$ jest izomorfizmem poza $p$, natomiast
jej włókno nad $p$ jest całym $\PP^1$.
Druga projekcja $q:B\to\PP^1$ ma włókna $\PP^1$.
Rzeczywiście, nad chartem $u\ne0$ piszemy $t=v/u$ i
\[
([\lambda:\lambda t:\mu],[1:t]),\qquad[\lambda:\mu]\in\PP^1.
\]
Nad chartem $v\ne0$ piszemy $s=u/v$ i otrzymujemy
$([s\lambda':\lambda':\mu'],[s:1])$.
Na przecięciu $s=1/t$ współrzędne włókien spełniają
$[\lambda':\mu']=[t\lambda:\mu]$. Ten wzór jest
projektywizacją przejścia wiązki
$\mathcal O_{\PP^1}(-1)\oplus\mathcal O_{\PP^1}$
w konwencji \emph{prostych} we włóknach: w lokalnych
ramkach $e_u=(1,t)$ i $e_v=(s,1)$ wiązki tautologicznej
zachodzi $e_u=t e_v$.
Dualizacja przekształca konwencję prostych w przyjętą
w definicji $\mathbb F_n$ konwencję jednowymiarowych ilorazów.
Dlatego
$B\cong\PP_{\PP^1}(\mathcal O\oplus\mathcal O(1))=\mathbb F_1$.
Wyjątkowa krzywa $E$ jest sekcją $q$, zadaną warunkiem
$\lambda=0$ ma $E^2=-1$; to jednocześnie
bezpośrednia identyfikacja powierzchni i jej sekcji,
bez odwoływania się do niewyjaśnionej jednoznaczności.
\end{przyklad}

\begin{uwaga}[Powierzchnie minimalne]
Gładką rzutową powierzchnię nazywamy \emph{minimalną},
jeżeli nie zawiera krzywych $(-1)$. Z twierdzenia
Castelnuova wynika, że każdą gładką rzutową powierzchnię
można, dopóki ma krzywe $(-1)$, próbować upraszczać
przez ich zdmuchiwanie; zakończenie tego procesu
uzasadnia np. zmiana liczby Picarda. Powierzchnia
$\mathbb F_1$ nie jest minimalna, natomiast
$\PP^2$ oraz $\mathbb F_n$ dla $n=0$ i $n\ge2$
są minimalne. Dla wymiernych powierzchni model minimalny
nie zawsze jest jednoznaczny, co odróżnia je od
odpowiedniego zagadnienia dla powierzchni o nieujemnym
wymiarze Kodairy.
\end{uwaga}

\section{Powierzchnie del Pezzo i antykanoniczne odwzorowania}
\label{sec:ga6-delpezzo}

\begin{definicja}[Powierzchnia del Pezzo i jej stopień]
Gładka rzutowa powierzchnia $S$ jest \emph{powierzchnią
del Pezzo}, jeśli $-K_S$ jest dywizorem obszernym
(ang. \emph{ample}). Przypomnijmy: wiązka liniowa $L$
jest obszerna, gdy dla pewnego $m>0$ liniowy układ
sekcji $L^{\otimes m}$ wyznacza zanurzenie zamknięte
$S\hookrightarrow\PP^N$. Jej \emph{stopień del Pezzo}
to liczba $d=K_S^2$.
\end{definicja}

\begin{przyklad}[Płaszczyzna i powierzchnia kwadrykowa]
Dla $\PP^2$ mamy $K_{\PP^2}\sim-3H$,
więc $(-K_{\PP^2})^2=9$ i $-K_{\PP^2}$ jest
obszerny. Zatem $\PP^2$ jest powierzchnią del Pezzo
stopnia dziewięć. Dla $\PP^1\times\PP^1$
mamy $-K=2F_1+2F_2$, co jest obszerne;
stopień wynosi osiem. Obie są powierzchniami
del Pezzo, mimo że nie są izomorficzne.
\end{przyklad}

\begin{stwierdzenie}[Stopień powierzchni otrzymanej z $\PP^2$]
\label{st:ga6-del-pezzo-degree}
Niech $\pi:S_r\to\PP^2$ będzie złożeniem
$r$ rozdmuchań w \emph{różnych} punktach
$p_1,\ldots,p_r$ płaszczyzny. Wówczas
\[
 K_{S_r}\sim-3H+E_1+\cdots+E_r,
 \qquad K_{S_r}^2=9-r.
\]
W sytuacji $r=6$ i punktów w położeniu ogólnym
(brak trzech współliniowych i brak stożkowej
przechodzącej przez wszystkie sześć punktów)
dywizor $-K_{S_6}=3H-\sum E_i$ jest nawet
\emph{bardzo obszerny}: jego układ zupełny
wyznacza zanurzenie $S_6\hookrightarrow\PP^3$
jako gładką powierzchnię kubiczną.
\end{stwierdzenie}
\begin{proof}
Formułę na $K_{S_r}$ otrzymujemy, stosując kolejno
stwierdzenie~\ref{st:ga6-Kblowup}. Ponieważ
$H^2=1$, $H\cdot E_i=0$, $E_i^2=-1$
oraz $E_i\cdot E_j=0$ dla $i\ne j$, dostajemy
$K^2=9-r$. Założenie o różnych punktach zapewnia,
że dzielniki wyjątkowe nie przecinają się.

Dla ostatniego zdania potrzeba osobnego twierdzenia
o \emph{bardzo obszerności} układu kubik
płaskich przechodzących przez sześć punktów
w położeniu ogólnym. Nie ukrywamy go pod formułą
$K^2=3$. Można jednak od razu wyjaśnić wymiar
układu: przestrzeń wszystkich jednorodnych form
stopnia trzy w trzech zmiennych ma wymiar
$\binom{3+2}{2}=10$. Sześć punktów nakłada
sześć niezależnych warunków liniowych: dla każdego
$p_i$ istnieje stożkowa przechodząca przez pięć
pozostałych, lecz nie przez $p_i$, a jej iloczyn
z prostą nieprzechodzącą przez $p_i$ rozdziela
wartość w $p_i$ od wartości w pozostałych punktach.
Otrzymujemy zatem czterowymiarową przestrzeń
sekcji $H^0(S_6,\mathcal O_{S_6}(-K))$.
Twierdzenie o bardzo obszerności tej konkretnej
wiązki zapewnia, że odpowiadający morfizm do
$\PP^3$ jest zanurzeniem, a nie tylko odwzorowaniem
wymiernym. Obraz ma stopień
$(-K)^2=3$, a więc jest powierzchnią kubiczną.
\end{proof}

\begin{uwaga}[Jaki wynik przyjęliśmy?]
Twierdzenie o bardzo obszerności $-K$ dla
sześciu punktów w położeniu ogólnym można udowodnić,
sprawdzając, że kubiki przez te punkty rozdzielają
punkty i kierunki styczne na rozdmuchaniu.
Pełny dowód wymagałby systematycznych narzędzi
do badania liniowych układów krzywych, więc
\emph{tę jedną część} przyjęliśmy jako wynik
zewnętrzny. Wszystkie dalsze obliczenia klas,
przecięć i liczby prostych wykonamy jawnie.
\end{uwaga}

\section{Dokładnie 27 prostych na gładkiej powierzchni kubicznej}
\label{sec:ga6-27}

\begin{twierdzenie}[Klasyczny wynik o 27 prostych]
\label{tw:ga6-27-lines}
Niech $k$ będzie algebraicznie domkniętym ciałem
charakterystyki zero, a
$S=\V_+(F)\subseteq\PP^3_k$ gładką powierzchnią
kubiczną, tzn. $F$ jest jednorodnym wielomianem
stopnia trzy. Wtedy $S$ zawiera dokładnie
$27$ prostych projektowych.
\end{twierdzenie}

Wyjaśnimy najpierw jedyny głęboki składnik dowodu.
Klasyczne twierdzenie o klasyfikacji powierzchni
del Pezzo stopnia trzy mówi, że \emph{każda} gładka
powierzchnia kubiczna nad $k$ powstaje jako
rozdmuchanie $\PP^2_k$ w sześciu punktach w położeniu
ogólnym, a jej zanurzenie kubiczne jest określone
pełnym układem antykanonicznym. Wynik ten przyjmujemy;
nie wynika on jedynie z istnienia $27$ krzywych
na pewnym konkretnym przykładzie.

\begin{proof}[Dowód twierdzenia~\ref{tw:ga6-27-lines}, przy użyciu powyższej klasyfikacji]
Niech $\pi:S\to\PP^2$ będzie rozdmuchaniem
sześciu punktów $p_1,\ldots,p_6$. Oznaczmy
przez $H$ pullback prostej, a przez $E_i$
krzywe wyjątkowe. Rozważmy antykanoniczny
dywizor $L=-K_S=3H-\sum_{i=1}^6E_i$.
Ponieważ $L$ definiuje zanurzenie $S\hookrightarrow\PP^3$,
krzywa całkowita $C\subseteq S$ jest prostą
w $\PP^3$ wtedy i tylko wtedy, gdy stopień
$\mathcal O_S(L)|_C$, czyli $L\cdot C$, wynosi jeden
(w kierunku odwrotnym korzystamy z faktu, że
nierozkładalna krzywa projektowa stopnia jeden
jest prostą).

\emph{Krok 1: każda prosta ma samoprzecięcie $-1$.}
Jeśli $C$ jest prostą, to $C\cong\PP^1$,
$(-K_S)\cdot C=1$. Ze wzoru adjunkcji
$-2=C^2+K_S\cdot C=C^2-1$, więc $C^2=-1$.

\emph{Krok 2: sześć krzywych wyjątkowych.}
Dla każdej $E_i$ mamy $E_i^2=-1$ i
$L\cdot E_i=1$. Ich obrazy są prostymi;
daje to sześć różnych prostych.

\emph{Krok 3: kandydaci niebędący krzywymi wyjątkowymi.}
Jeśli $C$ nie jest jedną z $E_i$, jej obraz przez
$\pi$ jest nierozkładalną krzywą płaską stopnia
$a\ge1$. Niech $b_i\ge0$ będą krotnościami
tej krzywej w punktach $p_i$. Wtedy
$[C]=aH-\sum_{i=1}^6b_iE_i$. Z warunków
$L\cdot C=1$ oraz $C^2=-1$ otrzymujemy układ
\[
 3a-\sum_{i=1}^6b_i=1,\qquad
 a^2-\sum_{i=1}^6b_i^2=-1.
\]
Czyli $\sum b_i=3a-1$ oraz
$\sum b_i^2=a^2+1$. Z nierówności Cauchy'ego--Schwarza
\[
 (3a-1)^2=\left(\sum b_i\right)^2
 \le6\sum b_i^2=6(a^2+1)
\]
wynika $3a^2-6a-5\le0$. Ponieważ $a\ge1$
jest liczbą całkowitą, musi zachodzić $a=1$ albo $a=2$.

\emph{Krok 4: przypadek $a=1$.}
Mamy $\sum b_i=2$ oraz $\sum b_i^2=2$.
Dla $b_i\in\Z_{\ge0}$ równość
$\sum b_i(b_i-1)=0$ wymusza, aby każde
$b_i$ było zerem albo jedynką. Dokładnie dwie
liczby $b_i$ są równe jeden. Jest zatem
$\binom{6}{2}=15$ możliwych klas
$H-E_i-E_j$, $i<j$. Każdą z nich reprezentuje
transformata ścisła jedynej prostej płaskiej
przechodzącej przez $p_i,p_j$. Założenie
ogólnego położenia wyklucza trzecie $p_l$
na tej prostej. Wszystkie $15$ są różne.

\emph{Krok 5: przypadek $a=2$.}
Teraz $\sum b_i=5$ i $\sum b_i^2=5$,
więc znowu pięć liczb $b_i$ wynosi jeden,
a szósta wynosi zero. Otrzymujemy sześć klas
$2H-\sum_{j\ne i}E_j$, $1\le i\le6$.
Przez pięć punktów w położeniu ogólnym przechodzi
dokładnie jedna stożkowa. Jest ona gładka:
stożkowa rozkładalna byłaby sumą dwóch prostych,
z których przynajmniej jedna zawierałaby trzy
z pięciu punktów. Nie przechodzi przez szósty
punkt na mocy ogólnego położenia. Jej transformata
ścisła ma podaną klasę i jest prostą w kubicznym
zanurzeniu.

\emph{Krok 6: kompletność.}
Każda prosta należy do jednej z kategorii z
kroków 2--5; każda podana krzywa istnieje,
ma przecięcie z $L$ równe jeden i jest inna
od pozostałych. Łącznie jest więc
$6+15+6=27$ prostych.
\end{proof}

\begin{center}
\begin{minipage}{\linewidth}
\centering
\begin{tikzpicture}[>=Stealth,every node/.style={font=\small}]
 \node[draw,rounded corners,align=center,text width=3.7cm,inner sep=9pt] (a) at (0,0)
 {\textbf{6} wyjątkowych\\$E_1,\ldots,E_6$};
 \node[draw,rounded corners,align=center,text width=3.7cm,inner sep=9pt] (b) at (5.0,0)
 {\textbf{15} prostych przez pary\\$H-E_i-E_j$};
 \node[draw,rounded corners,align=center,text width=3.9cm,inner sep=9pt] (c) at (10.1,0)
 {\textbf{6} stożkowych przez piątki\\$2H-\sum_{j\ne i}E_j$};
 \draw[->] (a.east)--(b.west);
 \draw[->] (b.east)--(c.west);
 \node[align=center] at (5.1,-1.3)
  {$6+15+6=27$ różnych krzywych $(-1)$\\obrazowanych na proste przez $|-K_S|$};
\end{tikzpicture}

{\small Rysunek 6.3. Trzy typy prostych na gładkiej powierzchni kubicznej. Strzałki porządkują klasy, a nie oznaczają odwzorowań między krzywymi.}\par
\end{minipage}
\end{center}

\begin{przyklad}[27 prostych na konkretnej powierzchni Fermata]
\label{prz:ga6-fermat}
Rozważmy powierzchnię
$S_F=\V_+(X_0^3+X_1^3+X_2^3+X_3^3)\subseteq\PP^3_\C$.
Jest gładka, bo wspólne znikanie pochodnych
$3X_0^2,\ldots,3X_3^2$ wymuszałoby znikanie
wszystkich współrzędnych, co w przestrzeni rzutowej
jest niemożliwe. Podzielmy zbiór indeksów
$\{0,1,2,3\}$ na dwie nieuporządkowane pary
$\{i,j\}\sqcup\{l,m\}$; są dokładnie trzy takie
podziały. Wybierzmy dowolne $\zeta,\eta\in\C$
spełniające $\zeta^3=\eta^3=-1$ i nałóżmy
równości liniowe $X_i=\zeta X_j$,
$X_l=\eta X_m$. Wyznaczają one prostą projektową,
a jej każdy punkt należy do $S_F$, ponieważ
$(\zeta^3+1)X_j^3+(\eta^3+1)X_m^3=0$.
Z różnych par współczynników w ustalonym podziale wynikają różne proste;
różne podziały również nie dają tej samej prostej, gdyż
punkty z dwiema zerowymi współrzędnymi należą do niej tylko
wtedy, gdy obie niezerowe współrzędne tworzą jedną z wybranych par.
Mamy $3$ podziały i $3\cdot3=9$ wyborów pary
$(\zeta,\eta)$, łącznie $27$ różnych prostych.
Poprzednie twierdzenie zapewnia, że nie ma już innych.
To przykład, w którym słynny wynik można najpierw
\emph{zobaczyć w równaniach}, a potem zrozumieć,
dlaczego musi zachodzić dla każdej gładkiej powierzchni
kubicznej, a nie tylko dla tej symetrycznej.
\end{przyklad}

\begin{uwaga}[Dlaczego położenie ogólne jest potrzebne?]
Jeśli trzy rozdmuchiwane punkty leżą na jednej prostej,
jej transformata ścisła ma klasę $H-E_1-E_2-E_3$
i kwadrat $1-3=-2$. Ponadto
$(-K)\cdot(H-E_1-E_2-E_3)=3-3=0$,
co wyklucza obszerność $-K$ na tej powierzchni
(z obszerności wynika dodatni stopień na każdej
krzywej całkowitej). W szczególności nie można
bez sprawdzenia ogólnego położenia stosować
wyliczenia $27$ prostych.
\end{uwaga}

\section{Powierzchnie K3: przykład o innej geometrii}
\label{sec:ga6-k3}

\begin{definicja}[Powierzchnia K3]
Gładką, spójną, rzutową powierzchnię $S$ nad $\C$
nazywamy \emph{powierzchnią K3}, gdy
$\omega_S\cong\mathcal O_S$ oraz
$H^1(S,\mathcal O_S)=0$. Tutaj
$\omega_S=\det\Omega_S^1$ jest wiązką kanoniczną,
a $H^1$ oznacza pierwszą kohomologię snopa strukturalnego.
\end{definicja}

\begin{twierdzenie}[Gładka powierzchnia stopnia czwartego jest K3]
\label{tw:ga6-quartic-k3}
Niech $S=\V_+(F)\subseteq\PP^3_\C$, gdzie $F$
jest jednorodnym wielomianem stopnia cztery, a $S$
jest gładka. Wtedy $S$ jest powierzchnią K3.
\end{twierdzenie}
\begin{proof}
Z adjunkcji dla gładkiej hiperpowierzchni
(udowadnianej tak samo przez ciąg konormalny jak
w twierdzeniu~\ref{tw:ga6-adjunction}, z wyznacznikami
wiązek odpowiedniego rzędu) wynika
$\omega_S\cong(\omega_{\PP^3}\otimes
\mathcal O_{\PP^3}(4))|_S$.
Ponieważ $\omega_{\PP^3}\cong\mathcal O_{\PP^3}(-4)$,
otrzymujemy $\omega_S\cong\mathcal O_S$.

Aby obliczyć pierwszą kohomologię strukturalną,
skorzystajmy z dokładnego ciągu snopów
\[
0\longrightarrow\mathcal O_{\PP^3}(-4)
\xrightarrow{\cdot F}\mathcal O_{\PP^3}
\longrightarrow\mathcal O_S\longrightarrow0.
\]
Odpowiedni fragment długiego ciągu kohomologii ma postać
$H^1(\PP^3,\mathcal O_{\PP^3})\to
H^1(S,\mathcal O_S)\to
H^2(\PP^3,\mathcal O_{\PP^3}(-4))$.
Skrajne grupy są zerowe: ogólne obliczenie kohomologii
wiązek $\mathcal O_{\PP^n}(m)$ daje
$H^i(\PP^n,\mathcal O(m))=0$ dla $0<i<n$,
niezależnie od całkowitego $m$. W szczególności
$H^1(S,\mathcal O_S)=0$.
Pozostałe założenia z definicji (spójność i
rzutowość) są spełnione: gładka hiperpowierzchnia
rzutowa dodatniego stopnia w $\PP^3$ jest spójna;
można to również zobaczyć z
$H^0(S,\mathcal O_S)=\C$ otrzymanego z powyższego
ciągu. Tym samym $S$ jest powierzchnią K3.
\end{proof}

\begin{przyklad}[K3 nie może zawierać krzywej $(-1)$]
\label{prz:ga6-k3-no-minus1}
Jeśli $C\cong\PP^1$ jest gładką krzywą na powierzchni K3,
to $K_S\sim0$, więc ze wzoru adjunkcji
$-2=C^2+K_S\cdot C=C^2$.
Każda gładka krzywa wymierna ma zatem
samoprzecięcie $-2$, a nie $-1$. Wynika stąd,
że powierzchnia K3 jest minimalna. Nie oznacza to,
że nie może zawierać żadnych gładkich krzywych
wymiernych: na przykład proste leżące na niektórych
gładkich kwartykach są krzywymi $(-2)$.
\end{przyklad}

\begin{stwierdzenie}[Genus przekroju hiperpłaszczyzną K3]
\label{st:ga6-quartic-plane-section}
Niech $S\subseteq\PP^3_\C$ będzie gładką kwartyką,
a $C=S\cap H$ jej gładkim przekrojem płaszczyzną
$H\cong\PP^2$. Wtedy $g(C)=3$.
\end{stwierdzenie}
\begin{proof}
Na powierzchni $S$ klasa takiego przekroju
jest oznaczana przez $h=\mathcal O_S(1)$.
Jej samoprzecięcie jest równe stopniowi $S$,
czyli $h^2=4$: przecięcie dwóch ogólnych
płaszczyzn z $S$ składa się z czterech punktów
liczonych z krotnościami (twierdzenie Bézouta).
Ponieważ $K_S\sim0$, adjunkcja daje
$2g(C)-2=(K_S+h)\cdot h=4$.
Stąd $g(C)=3$. Jest to zarazem wzór na rodzaj
gładkiej płaskiej krzywej stopnia cztery z rozdziału 4.
\end{proof}

\begin{uwaga}[Wymierna powierzchnia a K3]
Powierzchnia wymierna ma $H^0(S,\omega_S)=0$,
a K3 ma $H^0(S,\omega_S)\cong\C$.
Wymiar $p_g(S):=h^0(S,\omega_S)$ jest
niezmiennikiem biwymiernym gładkich powierzchni
rzutowych (jest to osobne twierdzenie, wynikające
z zachowania regularnych form dwuwymiarowych
pod rozdmuchaniami i rozszerzania przez zbiory
kowymiaru dwa). Zatem żadna powierzchnia K3
nie może być wymierna.
To nie jest tylko różnica w liczbie punktów:
dwie gładkie powierzchnie mogą być obie rzutowe
nad $\C$, a ich geometria biwymierna może być
całkowicie odmienna.
\end{uwaga}

\section{Niezmienniki powierzchni i mapa dalszej teorii}
\label{sec:ga6-invariants}

\begin{definicja}[Niezmienniki kohomologiczne]
Dla gładkiej spójnej rzutowej powierzchni $S$ nad
$\C$ definiujemy:
\[
 q(S)=\dim_{\C}H^1(S,\mathcal O_S),\qquad
 p_g(S)=\dim_{\C}H^0(S,\omega_S),\qquad
 \chi(\mathcal O_S)=\sum_{i=0}^2(-1)^i
 \dim_{\C}H^i(S,\mathcal O_S).
\]
Liczbę $q(S)$ nazywamy \emph{nieregularnością},
a $p_g(S)$ \emph{rodzajem geometrycznym} powierzchni.
Wszystkie występujące grupy są skończenie wymiarowe,
co wynika z twierdzenia o skończoności kohomologii
snopów koherentnych na schemacie właściwym.
\end{definicja}

\begin{stwierdzenie}[Wzór na charakterystykę holomorficzną]
\label{st:ga6-chi}
Dla opisanej powierzchni $S$ zachodzi
$\chi(\mathcal O_S)=1-q(S)+p_g(S)$.
\end{stwierdzenie}
\begin{proof}
Ponieważ $S$ jest spójna, właściwa i zredukowana
nad algebraicznie domkniętym ciałem, jej globalne
funkcje regularne są stałe, zatem
$h^0(S,\mathcal O_S)=1$. Z dualności Serre'a
dla gładkiej rzutowej powierzchni wynika
$H^2(S,\mathcal O_S)^\vee\cong H^0(S,\omega_S)$,
więc $h^2(S,\mathcal O_S)=p_g(S)$.
Podstawienie do definicji $\chi$ daje tezę.
\end{proof}

\begin{uwaga}[Klasy Cherna i wzór Noether]
Dwa klasyczne niezmienniki liczbowe powierzchni
to $K_S^2$ i stopień drugiej klasy Cherna
$c_2(S)=\deg c_2(T_S)$. Dla gładkiej rzutowej
powierzchni nad $\C$ \emph{wzór Noether} głosi
$12\chi(\mathcal O_S)=K_S^2+c_2(S)$.
Ponadto przez twierdzenie Gaussa--Bonneta--Cherna
$c_2(S)$ jest topologiczną charakterystyką Eulera
$S(\C)$. Na tym etapie oba wyniki są zewnętrzne.
Pełne definicje klas Cherna i dowody odpowiednich
formuł należą do rozdziału 7, a relacja z topologią
będzie tematem rozdziału 10. Nie korzystamy z nich
w dowodach o $27$ prostych ani o kwartykach K3.
\end{uwaga}

\begin{definicja}[Wymiar Kodairy]
Niech $S$ będzie gładką rzutową powierzchnią nad $\C$.
Dla $m\ge1$ oznaczmy
$P_m(S)=h^0(S,\omega_S^{\otimes m})$.
Jeżeli wszystkie $P_m(S)$ są zerowe, kładziemy
$\kappa(S)=-\infty$. W przeciwnym razie
$\kappa(S)$ jest największą liczbą całkowitą
$\kappa\in\{0,1,2\}$, dla której istnieje stała
$c>0$ i nieskończony podciąg dodatnich liczb
$m_j\to\infty$ taki, że
$P_{m_j}(S)\ge c\,m_j^{\kappa}$.
Fakt, że w tej definicji wystarczają wykładniki
$0,1,2$, a otrzymana liczba zgadza się z wymiarem obrazu
plurikanonicznego dla dostatecznie podzielnych $m$,
jest głębszym twierdzeniem (teoria wymiaru Iitaki),
które tutaj przyjmujemy.
Wymiar Kodairy mierzy zatem tempo wzrostu
liczby niezależnych sekcji potęg wiązki kanonicznej.
\end{definicja}

\begin{przyklad}[Wymiar Kodairy w dwóch skrajnych sytuacjach]
Dla $\PP^2$ wiązka kanoniczna to $\mathcal O(-3)$,
więc $H^0(\PP^2,\omega^{\otimes m})=0$
dla każdego $m>0$; stąd $\kappa(\PP^2)=-\infty$.
Dla K3 wiązka kanoniczna jest trywialna,
więc $P_m(S)=1$ dla każdego $m\ge1$.
Zatem $\kappa(S)=0$. To wyjaśnia, dlaczego
powierzchnie wymierne i K3 trafiają do różnych
części klasyfikacji Enriquesa--Kodairy.
\end{przyklad}

\begin{uwaga}[Jak odczytywać plan następnych rozdziałów?]
Rozdział 6 wykorzystał lokalne obliczenia rozdmuchań,
ale ograniczył ogólną teorię przecięć do postaci
podstawowych twierdzeń, z wyraźnym wskazaniem miejsc,
w których cytujemy wyniki zewnętrzne.
\begin{enumerate}[label=(\roman*)]
\item W rozdziale 7 zbudujemy cykle, grupy Chowa,
      teorię przecięć i klasy Cherna oraz wrócimy
      do dowodów wzorów liczbowych z tego rozdziału.
\item W rozdziale 8 omówimy geometrię biwymierną,
      modele minimalne, geometrię stożków krzywych
      i kontrakcje.
\item Rozdziały 9 i 10 połączą geometrię algebraiczną
      z analizą zespoloną i topologią: GAGA, Kähler,
      Hodge, Lefschetz i niezmienniki topologiczne.
\end{enumerate}
\end{uwaga}

% Literatura kontrolna i jawne zewnetrzne zaleznosci:
% Hartshorne, Algebraic Geometry, rozdzial V (surfaces).
% Beauville, Complex Algebraic Surfaces.
% Stacks Project, tag 0CBN: SNC versus NC, tag 01OF: blowups.
% M. Reid, Undergraduate Algebraic Geometry (cubic surfaces).
% R. Vakil, lectures on 27 lines on cubic surfaces.



% ===============================================================
% ROZDZIAL 7. TEORIA PRZECIEC I KLASY CHERNA — ITERACJA 9
% Zalozenie domyslne: k algebraicznie domkniete, char(k)=0.
% Wszystkie etykiety maja prefiks ga7:.
% ===============================================================
